% !TEX root = ../my-thesis.tex
%
\chapter{Foundations}
\label{sec:foundations}

This chapter introduces the formal and conceptual background the thesis builds
on. Section~\ref{sec:foundations:aaf} defines Dung's abstract argumentation
frameworks and the three semantics the pipeline computes.
Section~\ref{sec:foundations:typology} introduces the attack typology of Pollock
and ASPIC+ that the extraction step uses. Section~\ref{sec:foundations:am}
sketches argument mining, and Section~\ref{sec:foundations:llm} covers the
aspects of large language models relevant to this work.

\section{Abstract Argumentation Frameworks}
\label{sec:foundations:aaf}

Abstract argumentation frameworks (AFs), introduced by Phan Minh
Dung in 1995~\cite{dung1995}, are the most established formal model of
argumentation. The framework is deliberately abstract. It treats arguments as
atomic entities without internal structure and represents only the conflicts
between them.

\begin{definition}[Argumentation Framework]
An argumentation framework is a pair $AF = \langle \mathcal{A}, \mathcal{R}
\rangle$, where $\mathcal{A}$ is a finite set of arguments and $\mathcal{R}
\subseteq \mathcal{A} \times \mathcal{A}$ is a binary attack relation. We say
that an argument $a$ attacks an argument $b$ if $(a,b) \in \mathcal{R}$.
\end{definition}

The central question is which sets of arguments can rationally be accepted
together. The following notions, all due to Dung~\cite{dung1995}, build up to
an answer.

\begin{definition}[Conflict-freeness and defense]
A set $S \subseteq \mathcal{A}$ is \emph{conflict-free} if there are no
arguments $a, b \in S$ with $(a, b) \in \mathcal{R}$. An argument $a$ is
\emph{defended} by $S$ if for every $b$ with $(b, a) \in \mathcal{R}$ there is
some $c \in S$ with $(c, b) \in \mathcal{R}$.
\end{definition}

\begin{definition}[Admissibility and characteristic function]
A conflict-free set $S$ is \emph{admissible} if it defends every one of its
elements. The characteristic function $F_{AF} : 2^{\mathcal{A}} \to
2^{\mathcal{A}}$ maps a set to the arguments it defends,
$F_{AF}(S) = \{ a \in \mathcal{A} \mid S \text{ defends } a \}$.
\end{definition}

On top of admissibility, Dung defines several \emph{semantics}, each formalising
a different standard of collective acceptability. The following notion of a
complete extension is the common core of all semantics used in this thesis.

\begin{definition}[Complete extension]
An admissible set $S$ is a \emph{complete extension} if it contains every
argument it defends, that is, if $F_{AF}(S) = S$.
\end{definition}

A complete extension is a position that is both defensible and exhaustive. It
contains nothing it cannot defend, and it leaves out nothing it could defend.
This thesis computes three semantics that refine completeness, chosen to span a
wide spectrum of strictness.

\begin{itemize}
  \item The \textbf{grounded extension} is the least fixed point of $F_{AF}$,
    obtained by iterating $F_{AF}$ from the empty set. It is the smallest
    complete extension, it is the most cautious of the semantics, it always
    exists, and it is unique.
  \item A \textbf{preferred extension} is a maximal admissible set. Every
    preferred extension is complete, and it represents the largest internally
    coherent position a debate supports. A framework can have several preferred
    extensions.
  \item A \textbf{stable extension} is a conflict-free set $S$ that attacks every
    argument outside $S$. Stable extensions are the strictest of the three and
    are not guaranteed to exist.
\end{itemize}

The semantics form an inclusion hierarchy~\cite{dung1995}. Every stable
extension is preferred, every preferred extension is complete, and the grounded
extension is contained in every complete extension. The last property matters
twice in this thesis. Conceptually, the grounded extension is the uncontroversial
core that every rational position must accept. Computationally, it is the basis
of the solver optimization described in Section~\ref{sec:method:solver}.

The grounded semantics collapses a debate to a single cautious set, whereas
preferred and stable extensions, which may be multiple, reveal whether a topic
supports several internally consistent positions. This distinction is central to
the evaluation in Chapter~\ref{sec:evaluation}.

\subsection*{A worked example}

\begin{figure}[tb]
  \centering
  \begin{minipage}[c]{0.58\linewidth}
    \centering
    \begin{tikzpicture}[
      >={Stealth[length=2mm]},
      node distance=11mm and 13mm,
      arg/.style={draw, circle, minimum size=7mm, inner sep=1pt, font=\small},
      gr/.style={fill=gray!25},
    ]
      \node[arg]     (a) {$a$};
      \node[arg, right=of a] (b) {$b$};
      \node[arg, below=of a] (c) {$c$};
      \node[arg, right=of c, gr] (d) {$d$};
      \node[arg, left=of c, gr]  (e) {$e$};
      \draw[->] (a) to[bend left=15] (b);
      \draw[->] (b) to[bend left=15] (a);
      \draw[->] (e) -- (c);
      \draw[->] (c) -- (d);
    \end{tikzpicture}
  \end{minipage}\hfill
  \begin{minipage}[c]{0.38\linewidth}
    \centering
    \begin{tikzpicture}[
      >={Stealth[length=2mm]},
      arg/.style={draw, circle, minimum size=7mm, inner sep=1pt, font=\small},
    ]
      \node[arg] (x) at (90:1.1)  {$x$};
      \node[arg] (y) at (210:1.1) {$y$};
      \node[arg] (z) at (330:1.1) {$z$};
      \draw[->] (x) -- (y);
      \draw[->] (y) -- (z);
      \draw[->] (z) -- (x);
    \end{tikzpicture}
  \end{minipage}
  \caption{Two small argumentation frameworks. Left: the running example
  $AF_1$ with a mutual attack between $a$ and $b$. Filled nodes form the
  grounded extension $\{d, e\}$, and the two preferred extensions
  $\{a, d, e\}$ and $\{b, d, e\}$ extend it in incompatible ways. Right: the
  odd-length cycle $AF_2$, whose only complete extension is the empty set and
  which has no stable extension.}
  \label{fig:foundations:example}
\end{figure}

The difference between the three semantics is best seen on a small example.
Consider the framework $AF_1$ on the left of
Figure~\ref{fig:foundations:example}, with arguments
$\mathcal{A} = \{a, b, c, d, e\}$ and attacks
$\mathcal{R} = \{(a,b), (b,a), (e,c), (c,d)\}$. The structure mirrors a debate
in miniature. The arguments $a$ and $b$ are the cores of two opposing camps
that directly contradict each other, for instance the claim from the UKP
nuclear-energy topic that waste storage is safe and secure against the claim
that nuclear waste storage remains an unresolved safety problem. Each holds
only if the other fails, so the conflict runs both ways, which the mutual
attack between $a$ and $b$ captures. The argument $c$ raises an objection that
attacks a further pro claim $d$, for instance that nuclear power is virtually
carbon-free, and $c$ is in turn refuted by $e$.

Iterating the characteristic function from the empty set yields
$F_{AF}(\emptyset) = \{e\}$, since $e$ is the only unattacked argument, and
then $F_{AF}(\{e\}) = \{d, e\}$, since $e$ defends $d$ against its only
attacker $c$. A third step adds nothing, so the grounded extension is
$\{d, e\}$. It accepts the uncontroversial arguments and stays silent on the
conflict between $a$ and $b$. The preferred semantics goes further. Both
$\{a, d, e\}$ and $\{b, d, e\}$ are admissible, since $a$ and $b$ defend
themselves against each other, and neither set can be enlarged without a
conflict, so the framework has two preferred extensions, one per camp, each
containing the grounded extension as the inclusion hierarchy predicts. Both
sets are also stable, since each attacks every outside argument. The example
shows the division of labour that this thesis exploits. The grounded extension
identifies the safe core, and the preferred and stable extensions make the
incompatible ways of completing that core explicit.

The odd-length cycle $AF_2$ on the right of
Figure~\ref{fig:foundations:example}
illustrates why stable extensions are not guaranteed to exist. No single
argument in this example can defend itself against its attacker, so the empty set is the only
admissible and hence the only complete and preferred extension. The empty set
attacks nothing, so it is not stable, and $AF_2$ has no stable extension at
all.

\section{Attack Typology}
\label{sec:foundations:typology}

Dung's framework is abstract in that it does not say \emph{why} one argument
attacks another. For the extraction step, however, the kind of attack matters.
Pollock~\cite{pollock1992} distinguishes two kinds of defeaters in
natural-language reasoning, rebutting defeaters that contradict a conclusion
and undercutting defeaters that break the inference leading to it.
ASPIC+~\cite{modgil2014}, a structured argumentation framework that builds
Dung-style attack graphs from explicitly represented premises, inference
rules, and conclusions, adds undermining as an attack on a premise and gives
each attack a precise structural target. This yields the three-way typology
that the extraction step of this thesis uses.

\begin{itemize}
  \item A \textbf{rebuttal} contradicts the \emph{conclusion} of the target. An
    example is the pair ``nuclear power is clean'' against ``nuclear plants cause
    catastrophic contamination''.
  \item An \textbf{undercutting} attack challenges the \emph{inference} from
    premises to conclusion. The conclusion does not follow even if the premises
    hold. The claim ``nuclear power is safe because its CO\textsubscript{2}
    emissions are low'' is undercut by ``low emissions say nothing about
    accident or waste risks'', which grants the premise but severs the step to
    the conclusion.
  \item An \textbf{undermining} attack hits a \emph{premise} the target relies
    on, without engaging the conclusion directly. The claim ``containment makes
    plants safe'' is undermined by ``radiation cannot be contained
    indefinitely''.
\end{itemize}

A property that matters for the construction of the framework is that a rebuttal
is inherently mutual. If two conclusions negate the same proposition, then one is
true only if the other is false, so the conflict holds in both directions.
Without a preference or strength ordering, which Dung's abstract framework does
not provide, neither side takes precedence. The pipeline therefore symmetrizes
rebuttals, in line with the treatment of rebuttals in ASPIC+~\cite{modgil2014}.
Undercutting and undermining attacks, by contrast, remain directed. This
asymmetry has consequences for the cycles that appear in the resulting
frameworks, which Chapter~\ref{sec:evaluation} discusses.

\section{Argument Mining}
\label{sec:foundations:am}

Argument mining (AM) aims to automatically identify argumentative components,
such as claims and premises, and the relations between them from natural-language
text. Traditional approaches decompose the problem into subtasks including
argument detection, stance classification, and relation prediction, and address
them with supervised learning on annotated corpora~\cite{lawrence2019argument}.
More recently, LLMs have begun to replace task-specific classifiers with
prompt-based methods that cover these subtasks in a single
framework~\cite{li2025large}. Most existing work, however, evaluates AM
subtasks in isolation and stops before connecting the extracted relations to
formal reasoning, which leaves the question of which arguments are collectively
acceptable unanswered. This thesis targets exactly that gap.

\section{Large Language Models}
\label{sec:foundations:llm}

Large language models are neural networks trained on large text corpora to
predict the next token, and through that objective they acquire a broad ability
to follow instructions given in natural language. Two properties matter for
this thesis. Model behaviour is steered through \emph{prompts}, so the
wording of the task directly shapes the output, and even small changes to a
prompt can shift the results markedly~\cite{sclar2024}. This sensitivity is the very
lever that the typed extraction of this thesis relies on, and at the same time
a threat to stability, so the pipeline versions every prompt and
Section~\ref{sec:discussion:threats} discusses the observed effect of prompt
revisions. A sampling \emph{temperature} of zero in turn makes the output as
deterministic as the model allows, but not absolutely deterministic, since
server-side batching and numerical effects can still alter the sampled tokens.
